\documentclass[10pt,a4paper]{amsart}
\usepackage[margin=19mm]{geometry}
\usepackage{amsmath,amssymb,mathtools}
\usepackage[scr=rsfs,cal=euler]{mathalfa}
\usepackage{xfrac}
\usepackage[unicode,hidelinks,bookmarks=false]{hyperref}
\newcommand{\ie}{i.e.\ }
\newcommand{\cf}{cf.\ }
\newcommand{\resp}{respectively\ }
\newcommand{\Subsection}{\subsection}
\providecommand{\nobreakdash}{\nobreak\hbox{-}}
\setlength{\emergencystretch}{2em}
\multlinegap0pt
\usepackage{amssymb}
\usepackage{xcolor}
\usepackage[shortlabels]{enumitem}
\usepackage{cancel}
\newtheorem{lemma}{Lemma}
\renewcommand{\epsilon}{\varepsilon}
\renewcommand{\geq}{\geqslant}
\renewcommand{\leq}{\leqslant}
\newcommand{\tr}{\operatorname{tr}}
\title{Removability of non-isolated singularities for Einstein metrics and RCD spaces}
\author{Gioacchino Antonelli, G\'abor Sz\'ekelyhidi}
\date{Source: arXiv:2609.01464v1 (CC BY 4.0)}
\begin{document}
\maketitle

\noindent\emph{Source and licence.} Removability of non-isolated singularities for Einstein metrics and RCD spaces. Version arXiv:2609.01464v1 (CC BY 4.0), distributed by arXiv under the Creative Commons Attribution 4.0 International licence, http://creativecommons.org/licenses/by/4.0/, as recorded in the arXiv licence metadata for this exact version and shown on the abstract page https://arxiv.org/abs/2609.01464v1. Full licence notice: \texttt{licenses/arxiv-2609.01464-CC-BY-4.0.txt}.\par\smallskip

\noindent\textit{Editorial scope.} Counted unit: the article's lemma lem:Monotonicity (source0.tex lines 1662--1676). The article's complete original proof of it is delivered with it (source0.tex lines 1678--1712, one \texttt{proof} environment of 35 lines). No mathematics is written, repaired or re-derived: the statement and the proof are the article's own lines, in the article's own notation, and no retained line is substituted or re-flowed.  No further source-local context is needed: every cross-reference of every retained line resolves inside the two delivered spans.  Nothing was omitted from any retained span, so the package carries no omission record.  No retained line contains a citation, so the wrapper carries no bibliography and no bibliography entry.  No retained line contains a figure environment, an image-inclusion command or drawing code, so no image is delivered and the package declares no asset dependency.  Editorial formatting only, in addition: an AMS \texttt{amsart} preamble that restates the article's own declarations and macros used by the retained lines, namely the environments \texttt{lemma} on separate counters, together with the standard packages those lines need.  The retained environments print in this document as Lemma 1 in order of appearance, whereas the article prints the same items with the numbering of its own full document.  The complete native source of the article as distributed in the arXiv e-print for this exact version is bundled unchanged as \texttt{sources/209115/source0.tex}.
\begin{lemma}\label{lem:Monotonicity}
Let $n\geq2$, and let $\sigma$ be a locally finite, symmetric,
positive semidefinite matrix-valued Radon measure on $\mathbb R^n$. Put
$\mu:=\operatorname{tr}\sigma$ and $M(r):=\mu(B_r(0))$. If
\begin{equation}
 \operatorname{div}\left(\sigma-\frac12\mu I\right)=0,
 \label{eq:stationary-assumption}
\end{equation}
then
\begin{equation}
 \frac{M(s)}{s^{n-2}}\leq\frac{M(R)}{R^{n-2}}
 \qquad \forall 0<s<R.
 \label{eq:stationary-monotonicity}
\end{equation}
\end{lemma}

\begin{proof}
Let $\rho_\varepsilon$ be a standard nonnegative mollifier and set
$\sigma_\varepsilon:=\rho_\varepsilon*\sigma$ and
$\mu_\varepsilon:=\operatorname{tr}\sigma_\varepsilon$.
Convolution preserves positive semidefiniteness and
\eqref{eq:stationary-assumption}. Define
$
 T_\varepsilon:=\frac12\mu_\varepsilon I-
       \sigma_\varepsilon$.
Then $\operatorname{div}T_\varepsilon=0$ and
$\operatorname{tr}T_\varepsilon=(n-2)\mu_\varepsilon/2$. We apply the divergence
theorem to the vector field $T_\varepsilon x$ on $B_r$. Since $\operatorname{div} (T_\epsilon x) = \tr T_\epsilon$, this gives
\[
 \frac{n-2}{2}M_\varepsilon(r)
 =r\int_{\partial B_r}T_\varepsilon(e_r,e_r)
 =\frac r2M_\varepsilon'(r)
   -r\int_{\partial B_r}
       \sigma_\varepsilon(e_r,e_r),
\]
where $e_r$ is the unit radial vector field, and where $M_\varepsilon(r)=\mu_\epsilon(B_r(0))$. Thus
\begin{equation}
 rM_\varepsilon'(r)-(n-2)M_\varepsilon(r)
 =2r\int_{\partial B_r}
       \sigma_\varepsilon(e_r,e_r)\geq0.
 \label{eq:stationary-derivative}
\end{equation}
Equivalently, $
 \frac{\mathrm{d}}{\mathrm{d}r}\left(r^{2-n}M_\varepsilon(r)\right)\geq0$. 
It follows that \eqref{eq:stationary-monotonicity} holds with
$M_\varepsilon$ in place of $M$.
If $r$ is a continuity radius for $\mu$, weak convergence
$\mu_\varepsilon\rightharpoonup\mu$ gives
$M_\varepsilon(r)\to M(r)$. We may therefore pass to the limit first at
continuity radii and then at arbitrary radii by monotone approximation.
This proves \eqref{eq:stationary-monotonicity}. \end{proof}

\end{document}
